\documentclass[12pt,a4paper]{report}
\usepackage{graphicx}
\usepackage{longtable}
\usepackage{array}
\usepackage{hyperref}
\usepackage{mathptmx}
\usepackage{setspace}

\begin{document}
\doublespacing

\title{Game-Theory-Based Balance: Formalizing the Machine Unlearning Problem as a Game Between Model Provider and Data Provider}
\author{Abhishek Hupele \\ Research Scholar-240900465 \\ School of Computer Engineering, Manipal Institute of Technology, Manipal}
\date{March 2026}
\maketitle

\section*{Supervision}
\begin{itemize}
    \item Dr. Sanjay Singh, Professor (Guide), School of Computer Engineering, MIT Manipal
    \item Dr. P.C. Siddalingaswamy, Associate Dean (Co-Guide), School of Computer Engineering, MIT Manipal
    \item Dr. Naveen Kulkarni, Director – Product Engineering - AI, EdgeVerve Systems Limited
\end{itemize}
\newpage
\tableofcontents
\listoffigures
\listoftables

\chapter*{Glossary}
\begin{longtable}{>{\bfseries}p{3cm} p{10cm}}
AI & Artificial Intelligence \\
CNN & Convolution Neural Network \\
DNN & Deep Neural Network \\
LLM & Large Language Model \\
LSTM & Long Short Term Memory \\
MIA & Membership Inference Attack \\
ML & Machine Learning \\
MUL & Machine Unlearning \\
NPTEL & National Programme on Technology Enhanced Learning \\
PCA & Principal Component Analysis \\
PII & Personally Identifiable Information \\
RAI & Responsible AI \\
RNN & Recurrent Neural Network \\
SISA & Sharded, Isolated, Sliced, Aggregated \\
SVM & Support Vector Machine \\
\end{longtable}

\chapter{Introduction}
\section{Introduction to Machine Unlearning}
Machine Unlearning means making a machine forget specific training data so that it behaves as if the data was never used. Reasons include withdrawal of consent, unethical acquisition, adversarial data, copyright issues, or compliance requirements.  
Challenges: stochasticity, incrementality, catastrophe of unlearning.  
Main approaches:
\begin{itemize}
    \item Centralized Unlearning
    \item Federated (Distributed) Unlearning
    \item Unlearning Verification
    \item Privacy and Security Issues
\end{itemize}

\section{Introduction to Game Theory}
Game theory studies strategic interactions between players.  
Key concepts: players, strategies, payoffs, information, equilibrium (Nash Equilibrium).  
Applications span economics, political science, biology, computer science, business, and social sciences.

\section{Intersection of Game Theory and Privacy}
Game theory helps balance privacy and utility. In machine unlearning, over-unlearning harms performance, under-unlearning risks privacy leaks.  
This proposal formalizes unlearning as a game between model provider and data provider.

\section{Problems Identified}
\begin{itemize}
    \item PII misuse (Cambridge Analytica, Apple Card bias, Amazon hiring bias, Clearview AI, Target pregnancy prediction, Gartner’s GenAI misuse forecast).
    \item Need for balance between privacy and utility.
    \item Technical difficulty of certified removal.
    \item Lack of universal unlearning frameworks.
    \item Performance degradation and leakage risks.
    \item Variability in models, modalities, adversaries.
\end{itemize}

\section{Proposed Solutions}
Formalize machine unlearning as a game-theoretic problem.  
Develop generalized frameworks applicable across ML models, modalities, and adversary settings.  
Flexible scope: start with classical ML and neural networks, expand to transformers and federated settings.

\chapter{Literature Review}
\begin{longtable}{|c|p{4cm}|p{3cm}|c|p{3cm}|p{3cm}|p{3cm}|}
\hline
Sl No & Title & Authors & Year & Methodology & Result & Limitations \\
\hline
1 & Responsible AI: Principles, Practices, and Applications & S. Wesselby & 2025 & Online course & Practitioner guidance & Not peer-reviewed \\
\hline
2 & Machine Unlearning: A Comprehensive Survey & Z. Wang et al. & 2024 & Survey taxonomy & Overview of methods & Limited empirical validation \\
\hline
3 & Machine Unlearning: A Survey & H. Xu et al. & 2023 & ACM Computing Surveys & Establishes unlearning categories & Limited streaming data coverage \\
\hline
4 & MLPro-GT—game theory and dynamic games modeling in Python & S. Yuwono et al. & 2024 & Framework paper & Open-source GT library & Limited templates \\
\hline
% Continue table entries as needed...
\end{longtable}

\chapter{Research Gaps}
% Content continues...

\chapter{Objectives}
% Content continues...

\chapter{Methodology}
% Content continues...

\chapter{Expected Outcome}
% Content continues...

\chapter{Importance of Proposed Research and SDG Link}
% Content continues...

\chapter{Research Time Plan}
% Content continues...

\chapter{Pilot Study}
% Content continues...

\chapter{Details of Expenses and Source of Funding}
% Content continues...

\chapter{Course Work Details}
% Content continues...

\chapter{References}
% Insert bibliography here

\end{document}